% Section 5.1, from lectures/L05/appendix/a_uart_regs.md.

\section{Designing \code{uart_regs.vhd}}
\label{c5:sec:uart_regs}\label{c5:app:a}
The datapath speaks in pulses and levels: a \code{valid} that is high for one clock, a \code{busy}
that rises and falls, a \code{tick}. Software speaks in \textbf{registers}: it writes a byte to a
fixed address, polls a status word, reads a byte back. The register bank is the translator between
the two, and it is the layer the whole protocol (Appendix~\ref{proto:ch}) describes.

It owns two of the \code{fifo} you built in \chapref{4} (\secref{c4:app:a}), the \code{CTRL},
\code{BAUD_DIV} and \code{ERROR_FLAGS} registers, and it computes the read-only \code{STATUS} word.
It is reached over a small register bus that the provided SPI bridge drives; the bridge, not this
module, is what turns SPI transactions into \code{reg_write} strobes, so \code{uart_regs} can trust
that a \code{reg_write} it sees is a real, completed write.

\subsection{Interface}
\label{c5:sec:uart_regs_interface}

\bookfigure{uart_regs}{Module \code{uart_regs}}{c5:fig:uart_regs}

\begin{booktable}{@{}p{7.6em}p{2em}p{15.2em}L@{}}
  \thead{Port} & \thead{Dir} & \thead{Type} & \thead{Meaning} \\
  \midrule
  \code{clock}, \code{reset_s2_n} & in & \code{std_logic} & Clock, and an active-low synchronized
    reset. \\
  \code{reg_addr} & in & \code{std_logic_vector(3 downto 0)} & Register index
    (\code{offset / 4}). \\
  \code{reg_wdata} & in & \code{std_logic_vector(31 downto 0)} & Write data. \\
  \code{reg_write} & in & \code{std_logic} & One-cycle commit strobe from the bridge. \\
  \code{tx_pop} & in & \code{std_logic} & The transmitter has taken the front byte. \\
  \code{rx_byte} & in & \code{std_logic_vector(7 downto 0)} & A received byte, from the receiver. \\
  \code{rx_push} & in & \code{std_logic} & Push \code{rx_byte} (the receiver's \code{valid}). \\
  \code{tx_busy} & in & \code{std_logic} & The transmitter is mid-frame. \\
  \code{frame_err} & in & \code{std_logic} & The receiver saw a bad stop bit. \\
  \code{reg_rdata} & out & \code{std_logic_vector(31 downto 0)} & Combinational read value for
    \code{reg_addr}. \\
  \code{baud_div} & out & \code{std_logic_vector(15 downto 0)} & \code{BAUD_DIV}, for
    \code{baud_gen}. \\
  \code{tx_byte} & out & \code{std_logic_vector(7 downto 0)} & TX FIFO front, to the transmitter. \\
  \code{tx_empty} & out & \code{std_logic} & TX FIFO empty. \\
  \code{rx_full} & out & \code{std_logic} & RX FIFO full. \\
\end{booktable}

The ports are declared all inputs first, then all outputs; \code{uart_regs_tb} binds them
positionally, so the order above is the contract. The testbench drives the register bus directly,
standing in for the bridge, and models the datapath side with plain signals, so the register
semantics are tested with no SPI in sight.

The register map itself (indices, offsets, bit positions) lives in the provided \code{uart_def}
package, so the bank and the testbench name registers and bits (\code{REG_STATUS},
\code{ST_TX_READY}, \ldots) instead of magic numbers.

\subsection{Register semantics (the whole point of this module)}
\label{c5:sec:uart_regs_semantics}
The map promises sticky, poll-able bits and FIFO-backed data. The datapath offers pulses and
levels. Bridging the two, bit by bit, is what the bank does. The three flag registers, from
\secref{c1:app:a}, are shown again in Figure~\ref{c5:fig:reg_fields}.

\bookfigure{reg_fields}{The bit fields of \code{STATUS}, \code{CTRL} and \code{ERROR_FLAGS}}
  {c5:fig:reg_fields}

\textbf{\code{STATUS} is computed, never stored.} Bit 0 (TX ready) is \code{not tx_full}, bit 1 (RX
valid) is \code{not rx_empty}, bit 2 (Error) is the OR of the error flags, and bit 3 (TX idle) is
\code{tx_empty and not tx_busy}. Software polls these; nothing writes them.

\textbf{\code{TX_DATA} is a write-triggered action, not storage.} A \code{reg_write} to its index
pushes \code{reg_wdata(7 downto 0)} into the TX FIFO. It is an edge event, one write for one byte,
and a separate feeder in \code{uart_top} (\secref{c1:app:b}) drains the FIFO into the transmitter.

\textbf{\code{RX_DATA} is a pure read; \code{RX_POP} is the advance.} Reading \code{RX_DATA} returns
the FIFO's front byte and changes nothing, so it obeys the same latch-once rule as \code{STATUS},
and discarding that byte is a \emph{separate} write to \code{RX_POP}. The split is deliberate. A read
that also popped the FIFO would be a read with a side effect, and it would need the same
commit-on-completion, abort-safe handling as a write: what should happen if the SPI transaction is
aborted mid-read? Keeping the read pure and the pop an explicit write sidesteps the question
entirely, and it is the single most important idea in this module.

\textbf{\code{ERROR_FLAGS} latches.} A \code{frame_err} pulse from the receiver sets the framing
bit, and it stays set until software writes \code{0} to clear it. \code{STATUS} bit 2 is just the OR
of these flags.

\textbf{\code{CTRL} and \code{BAUD_DIV} are plain read/write registers.} \code{baud_div} is driven
straight out to \code{baud_gen}, so a \code{BAUD_DIV} write takes effect on the next baud tick the
generator produces.

\code{CTRL} is different, and the difference matters: in this course it is \textbf{storage only, and
gates nothing}. Look at the port list above and you will not find a \code{ctrl} output, so there is
no path by which \code{CT_ENABLE}, the parity select, the stop-bit select or the two IRQ masks could
reach the datapath. They are stored, they read back, and nothing acts on them. Do \textbf{not} gate
your TX push or your receiver on \code{CT_ENABLE}: \code{uart_regs_tb} writes \code{CTRL} only in its last
case, long after its transmit case, and the system testbench \code{uart_top_tb} never writes it at
all, so a bank that waits to be enabled simply never transmits, and fails with messages that say
nothing about the cause: \code{TX FIFO front should be 0x5A, got 0x00!} in the one, a timeout
waiting for RX valid in the other. The bits exist in
\code{uart_def} because the register map is the shared contract with the C++ side and with the
protocol; wiring them to something real is a natural extension, not part of the build here.

The bank never worries about aborted or partial transactions. The bridge only asserts
\code{reg_write} on a completed, non-aborted write (Part~3 of the protocol, \secref{proto:part3}), so
a strobe here is always a real commit. That contract is exactly what lets \code{TX_DATA} and
\code{RX_POP} be simple single-cycle actions.

\subsection{Behaviour}
\label{c5:sec:uart_regs_behaviour}
Very little of this module is clocked. It owns two \code{fifo} instances, three stored registers,
and a pile of decode; the register bus itself is combinational, and the only state that advances on
an edge is \code{ctrl_reg}, \code{baud_div_reg}, \code{err_flags}, and whatever the two FIFOs do
internally.

Everything below is in the order you would type it: what to declare, what to instantiate, the
concurrent lines, then the one process. The blocks and what connects them come first, in
Figure~\ref{c5:fig:regs_internals}.

\bookfigure{regs_internals}{Inside \code{uart_regs}}{c5:fig:regs_internals}

Read the two queues against each other and the asymmetry is the whole module. The TX FIFO takes its
data from the bus and its pop from the datapath; the RX FIFO takes its data from the datapath and
its pop from the bus. Everything else is the decode that turns one \code{reg_write} into one action,
the \code{flags} the \code{STATUS} word is assembled from, and the mux that answers reads.

Then the control flow. The left half of the module is one clocked decode, the right half is pure
combinational read-back, and they never touch:

\begin{console}
on each rising edge, if reg_write = '1', decode reg_idx:
    REG_CTRL       : ctrl_reg(CT_TX_IRQ downto 0) <= reg_wdata(CT_TX_IRQ downto 0)
    REG_BAUD_DIV   : baud_div_reg <= reg_wdata(15 downto 0)
    REG_ERR_FLAGS  : err_flags    <= reg_wdata(2 downto 0)
    REG_TX_DATA    : nothing here - tx_push already pushed the FIFO
    REG_RX_POP     : nothing here - rx_pop already advanced the FIFO
    anything else  : ignored
then, last in the process, so that it beats a clear in the same cycle:
    if frame_err = '1' then err_flags(ER_FRAMING) <= '1'

combinationally, every cycle, reg_rdata follows reg_idx:
    reg_rdata <= (others => '0')                  -- default, then override
    when REG_STATUS    => reg_rdata              <= status_reg
    when REG_CTRL      => reg_rdata              <= ctrl_reg
    when REG_BAUD_DIV  => reg_rdata(15 downto 0) <= baud_div_reg
    when REG_RX_DATA   => reg_rdata(7 downto 0)  <= rx_front  -- never a pop
    when REG_ERR_FLAGS => reg_rdata(2 downto 0)  <= err_flags
    when others        => null                    -- TX_DATA, RX_POP, 7-15
\end{console}

A \code{TX_DATA} write is the one part of this that is easier to see in time than in structure. The
bridge holds the address and the word, raises \code{reg_write} for exactly one cycle, and the byte
lands on the edge inside it, as Figure~\ref{c5:fig:bus_timing} shows.

\bookfigure{bus_timing}{A \code{TX_DATA} write, cycle by cycle}{c5:fig:bus_timing}

\code{tx_push} is combinational, so it is high for the same cycle as \code{reg_write}; the FIFO
registers the push on the edge at the end of that cycle, which is where \code{tx_empty} falls and
\code{tx_byte} starts showing the byte. That one-cycle strobe is the shape \code{uart_regs_tb}'s
\code{bus_write} procedure drives, and it is why \code{TX_DATA} needs no storage of its own.

\paragraph{What to declare} The port list already gives you every input and output, so the
architecture adds only what the ports cannot hold:
\begin{itemize}
  \item \code{ctrl_reg[31:0]}: the control register (\code{CTRL}). Only bits 5--0 are ever written;
    bits 31--6 are reserved, so they stay at their reset value and read back as zero, exactly as
    \code{BAUD_DIV} and \code{ERROR_FLAGS} do with theirs.
  \item \code{baud_div_reg[15:0]}: the baud divider register (\code{BAUD_DIV}).
  \item \code{err_flags[2:0]}: the three error flags.
  \item \code{tx_push}, \code{rx_pop}: the two decoded write strobes.
  \item \code{tx_empty_s}, \code{tx_full_s}, \code{rx_empty_s}, \code{rx_full_s}, all
    \code{std_logic}: the four FIFO flags.
  \item \code{rx_front[7:0]}: the RX FIFO's front byte, which the read mux needs.
  \item \code{status_reg[31:0]}: the assembled \code{STATUS} value.
\end{itemize}

\paragraph{The two \code{fifo} instances} Both are \code{generic map(8, 8)}, and both bind
positionally:
\begin{itemize}
  \item The TX FIFO, in port order: \code{clock}, \code{reset_s2_n}, \code{reg_wdata(7 downto 0)},
    \code{tx_push}, \code{tx_pop}, \code{tx_byte}, \code{tx_empty_s}, \code{tx_full_s}.
  \item The RX FIFO, in port order: \code{clock}, \code{reset_s2_n}, \code{rx_byte},
    \code{rx_push}, \code{rx_pop}, \code{rx_front}, \code{rx_empty_s}, \code{rx_full_s}.
\end{itemize}
Read those two lines against each other and the asymmetry is the whole module: the TX queue takes
its data from the bus and its pop from the datapath, and the RX queue takes its data from the
datapath and its pop from the bus. \code{tx_byte} binds straight to the output port because nothing
in here reads it; \code{rx_front} is internal because the read mux does.

\paragraph{The concurrent assignments} One line each, outside any process. Everything the read mux
and the clocked process below need is built here:
\begin{itemize}
  \item \code{reg_idx} is \code{to_integer(unsigned(reg_addr))}, which needs
    \code{ieee.numeric_std}. Four address bits means 0 to 15, so the decode's \code{others} branch
    covers the reserved indices 7 to 15 for free.
  \item \code{tx_push} is \code{'1'} when \code{reg_write = '1'} and \code{reg_idx = REG_TX_DATA},
    else \code{'0'}.
  \item \code{rx_pop} is \code{'1'} when \code{reg_write = '1'} and \code{reg_idx = REG_RX_POP},
    else \code{'0'}.
  \item \code{tx_empty} follows \code{tx_empty_s}, and \code{rx_full} follows \code{rx_full_s}: the
    two flags that leave the module.
  \item \code{baud_div} follows \code{baud_div_reg}, straight out to the conversion in
    \code{uart_top}.
  \item \code{status_reg} is built one bit at a time, and every bit not named below is \code{'0'}:
  \begin{itemize}
    \item \code{status_reg(ST_TX_READY)} is \code{not tx_full_s}: there is room to push.
    \item \code{status_reg(ST_RX_VALID)} is \code{not rx_empty_s}: there is a byte to read.
    \item \code{status_reg(ST_ERROR)} is \code{'1'} when \code{err_flags} is not all zeros: the OR
      of the three.
    \item \code{status_reg(ST_TX_IDLE)} is \code{tx_empty_s and not tx_busy}: nothing queued
      \emph{and} no frame in flight. TX ready is about room, TX idle is about quiet; software waits
      on the second before reconfiguring.
    \item \code{status_reg(31 downto 4)} is \code{(others => '0')}, so the unused bits read back as
      the protocol promises. Every bit needs exactly one driver, which is why this line exists at
      all.
  \end{itemize}
\end{itemize}

\paragraph{The read mux} A combinational process: no clock, no register, one \code{case} on
\code{reg_idx}. Assign the zeros first and let each index override the bits it owns:

\begin{vhdlcode}
process(reg_idx, status_reg, ctrl_reg, baud_div_reg, rx_front, err_flags) is
begin
    reg_rdata <= (others => '0');            -- TX_DATA, RX_POP and 7-15 read zero
    case reg_idx is
        when REG_STATUS    => reg_rdata              <= status_reg;
        when REG_CTRL      => reg_rdata              <= ctrl_reg;
        when REG_BAUD_DIV  => reg_rdata(15 downto 0) <= baud_div_reg;
        when REG_RX_DATA   => reg_rdata(7 downto 0)  <= rx_front;  -- never pops
        when REG_ERR_FLAGS => reg_rdata(2 downto 0)  <= err_flags;
        when others        => null;
    end case;
end process;
\end{vhdlcode}

\noindent Four things that shape holds down:
\begin{itemize}
  \item \textbf{The default assignment comes first}, so every path through the process assigns
    \code{reg_rdata} and no latch is inferred. It also spares you five zero-extension expressions:
    assign the narrow register to the low bits and the leading zeros are already there.
  \item \textbf{\code{when others => null;} is honest here}, not lazy, precisely because of that
    default. Written without it, the same process would need \code{(others => '0')} spelled out in
    three more branches.
  \item \textbf{The reserved indices come free.} \code{reg_idx} is \code{natural range 0 to 15}, so
    \code{when others} covers 7 to 15 as well as the two write-only registers, which is what Part~2
    of the protocol asks for.
  \item \textbf{The sensitivity list is spelled out.} VHDL-93 has no \code{process(all)} and this
    course analyzes with \code{--std=93}, so every signal the process reads has to be listed, or
    simulation and synthesis disagree about a mux that looked fine on the page. The provided
    \code{spi_reg_bridge} does exactly this in \file{hw/spi_reg_bridge.vhd}, and is worth reading as
    the pattern.
\end{itemize}

Because that path is combinational, \code{reg_rdata} follows \code{reg_addr} within the cycle.
Nothing registers it on the way out: the bridge latches the word once, just after the command byte,
and \code{uart_regs_tb} allows it a nanosecond (\code{SETTLE_NS}) to settle before it samples.

\paragraph{The clocked process} It is \code{process(clock, reset_s2_n)}, with its reset asserted
asynchronously like every other module here, and it is the only \emph{clocked} process in the file
--- the read mux above is the other one, and it has no clock at all:
\begin{itemize}
  \item On reset (\code{reset_s2_n} is low):
  \begin{itemize}
    \item \code{ctrl_reg} is cleared, so every \code{CTRL} bit reads back low.
    \item \code{baud_div_reg} is cleared, so \code{BAUD_DIV} reads 0. A zero divider is not a legal
      input to \code{baud_gen}, which is why \code{uart_top} guards the conversion
      (\secref{c1:app:b}); the driver is expected to write a real divider before using the line.
    \item \code{err_flags} is cleared, so \code{STATUS} bit 2 reads low.
    \item Nothing else. The FIFOs take \code{reset_s2_n} on their own ports and come out empty,
      which is what makes the first \code{STATUS} read after reset show TX ready and TX idle set
      with RX valid clear.
  \end{itemize}
  \item On each rising edge of \code{clock}, while not in reset:
  \begin{itemize}
    \item First, if \code{reg_write} is high, decode \code{reg_idx}. \textbf{\code{REG_CTRL}}:
      \code{ctrl_reg(CT_TX_IRQ downto 0)} takes \code{reg_wdata(CT_TX_IRQ downto 0)}, so the six
      defined bits are stored and the reserved ones above them stay zero. Nothing else happens: no
      \code{ctrl} output exists, so the bits are stored, read back, and act on nothing.
      \textbf{\code{REG_BAUD_DIV}}: \code{baud_div_reg} takes \code{reg_wdata(15 downto 0)}. The new
      divider reaches \code{baud_gen} combinationally, so it applies from the next tick the
      generator produces. \textbf{\code{REG_ERR_FLAGS}}: \code{err_flags} takes
      \code{reg_wdata(2 downto 0)}, so the documented write of \code{0x0} clears all three latches.
      \textbf{\code{REG_TX_DATA}} and \textbf{\code{REG_RX_POP}}: nothing happens \emph{here}. Their
      whole effect is the concurrent strobe into a FIFO, which that FIFO registers on this same
      edge. Writing FIFO logic in this process as well is the most common way to get two bytes out
      of one write. \textbf{\code{REG_STATUS}}, \textbf{\code{REG_RX_DATA}}, and indices 7 to 15:
      ignored. A write to a read-only or reserved register does nothing at all.
    \item Then, \emph{after} that decode, if \code{frame_err} is high, set
      \code{err_flags(ER_FRAMING)}.
    \item Nothing else is clocked. No byte passes through a register on its way to
      \code{reg_rdata}, and no output is registered.
  \end{itemize}
\end{itemize}

\paragraph{Three things that catch people out}
\begin{itemize}
  \item \textbf{The order of the last two steps.} An \code{ERROR_FLAGS} clear and a fresh
    \code{frame_err} can land in the same cycle, and the last assignment in a process wins --- the
    same rule \code{uart_rx} leans on for its tick counter. Putting the latch after the write decode
    means the new error survives the clear; the other order loses a real error to a write issued
    before it happened, and a dropped error flag is far worse than a stale one.
  \item \textbf{The FIFOs already guard themselves.} A \code{TX_DATA} write while the TX FIFO is
    full is dropped by the FIFO, and an \code{RX_POP} on an empty RX FIFO does nothing, both because
    \chapref{4} guarded them on \code{not full} and \code{not empty}. The bank adds no guards of its
    own; \code{STATUS} is what the driver is supposed to poll first.
  \item \textbf{Nothing is gated on \code{CT_ENABLE}.} As above: no testbench writes \code{CTRL} before it transmits, so a bank that waits to be enabled never transmits.
\end{itemize}

Two things are deliberately absent. \code{ER_PARITY} and \code{ER_OVERRUN} have no producer in this
build, so they hold whatever a write last left in them, which is zero; \code{rx_full} is exported
precisely so that overrun detection can be added later without touching the map. And there is no TX
feeder in here: \code{tx_pop} is an \emph{input}, driven by the feeder in \code{uart_top}, so this
module never decides when the transmitter takes a byte.

\subsection{What the testbench pins down}
\label{c5:sec:uart_regs_tb}
\code{uart_regs_tb} works entirely over the register bus. It checks that \code{BAUD_DIV} writes and
reads back and drives \code{baud_div}, and that \code{STATUS} after reset shows TX ready and TX idle
set with RX valid and Error clear. On the transmit side, a \code{TX_DATA} write must appear at the
FIFO front and clear TX idle, and a \code{tx_pop} must empty it again. On the receive side, an
\code{rx_push} must set RX valid, \code{RX_DATA} must return the byte, a bare \code{RX_DATA} read
must \textbf{not} pop, and an \code{RX_POP} write must clear RX valid. Finally, a \code{frame_err}
must latch into \code{ERROR_FLAGS} and \code{STATUS}, and a write of \code{0} must clear it. Two
checks close the map off: a \code{CTRL} write of all ones must read back as \code{0x0000003F},
proving that the six defined bits are stored and the reserved ones above them are not, and a write
to reserved index 7 must be ignored, with the read returning zero.

The read-does-not-pop check is the one that pins down the read/pop split; get it wrong and RX valid
would clear on a plain read.

\subsection{Where it fits}
\label{c5:sec:uart_regs_fits}
\code{uart_regs} is the middle of \code{uart_top} (\secref{c1:app:b}): the bridge on one side
driving the register bus, the datapath (\code{baud_gen}, \code{uart_tx}, \code{uart_rx}) on the
other. Its two FIFOs are the buffers that let software and the serial line run at their own speeds.

\subsection{What's ahead}
The exercises in \secref{c5:sec:exercises} build \code{uart_regs}, then instantiate the bank into the
\code{uart_top} skeleton from \chapref{1}, the last block it was waiting for, and run the full
system testbench for the first time.
